% ------------------------------------------------------------------
%  Capitolo 8 — Ricordare per ricordare
%  Parte III
%  Ricerca di riferimento: ricerca 01 §5.1; 03 §2.5, §4.4; 05 §7
%  Voci bibliografiche principali: roediger2006, karpicke2008, karpicke2011, rowland2014, adesope2017, yang2021, butler2010, agarwal2019
%  Epigrafe: ✔ verificata sul testo (vedi ricerca 03 e registro, ricerca 00)
%  Scheda del capitolo: ricerca/06_indice-ragionato.md, capitolo 8
% ------------------------------------------------------------------
\chapter{Ricordare per ricordare}
\label{cap:recupero}

\epigrafe[gli esercizi conservano la memoria facendo ricordare di nuovo]{αἱ μελέται τὴν μνήμην σῴζουσι τῷ ἐπαναμιμνήσκειν}{Aristotele, De memoria et reminiscentia 451a12}

\iniziale{Q}{uesto capitolo} \segnaposto{apertura: da scrivere — vedi la scheda nell'indice ragionato}.

\section{Da Aristotele a Roediger}

\segnaposto{da scrivere}

\section{Gli esperimenti}

\segnaposto{da scrivere}

\section{Perché funziona}

\segnaposto{da scrivere}

\section{Come si fa}

\segnaposto{da scrivere}

\section{Gli errori più comuni}

\segnaposto{da scrivere}

\begin{daricordare}
\segnaposto{il principio del capitolo in due righe}
\end{daricordare}

\begin{domande}
  \item \segnaposto{domanda di recupero su questo capitolo}
  \item \segnaposto{domanda di applicazione a una materia che studi}
  \item \segnaposto{domanda di ripresa su un capitolo precedente}
\end{domande}

\begin{sintesi}
  \item \segnaposto{punto di sintesi}
\end{sintesi}

\finecapitolo
